\documentclass[a4paper]{ltjsarticle}

\usepackage{amsmath,amssymb,amsthm}
\usepackage{mathtools}
\usepackage{physics}
\usepackage{tikz-cd}
\usepackage{hyperref}
\usepackage{enumitem}

\hypersetup{
  colorlinks=true,
  linkcolor=blue,
  urlcolor=teal,
  citecolor=magenta
}

\title{離散時間フィルトレーションと測度論的構造塔の入門}
\author{CODEX (OpenAI)\thanks{TeXファイルおよび Lean ファイルを CODEX (OpenAI) が生成}}
\date{\today}

\newtheorem{definition}{定義}
\newtheorem{theorem}{定理}
\newtheorem{proposition}{命題}
\newtheorem{example}{例}
\theoremstyle{remark}
\newtheorem{remark}{補足}

\begin{document}

\maketitle

\begin{abstract}
本ノートは、Nicolas Bourbaki に着想を得た「構造塔」の言語を通して、確率論における離散時間フィルトレーションを学部生向けに解説する。特に、Lean ファイル \texttt{Claude\_Probability.lean} で形式化された \texttt{DiscreteFiltration} と \texttt{MeasurableTowerWithMin} の関係、さらには直積構成と簡単な例を紹介する。
\end{abstract}

\tableofcontents

\section{序論}

従来の確率論では、離散時間フィルトレーション $(\mathcal{F}_n)_{n\in \mathbb{N}}$ は、時間が進むにつれて情報が増える仕組みを表現する。本ノートでは、Bourbaki の「構造塔」の発想を活かし、フィルトレーションを \emph{層} の列として扱う。Lean 実装では、各層が集合の族ではなく可測空間 (\texttt{MeasurableSpace}) として格納され、最小層を与える関数 \texttt{minLayer} を通して停止時刻の考え方と自然に結びつく。

以下で扱う主要な Lean 中の定義は次の通りである。
\begin{itemize}[nosep]
  \item \texttt{DiscreteFiltration}：型 $\Omega$ 上の自然数添字付き可測空間列。
  \item \texttt{toMeasurableTower}：フィルトレーションを \texttt{MeasurableTowerWithMin} に変換。
  \item \texttt{product}：二つの離散フィルトレーションの直積。
\end{itemize}

\section{離散時間フィルトレーション}

\begin{definition}[離散時間フィルトレーション]
型 $\Omega$ に可測空間構造が与えられているとき、\texttt{Claude\_Probability.lean} の定義 \texttt{DiscreteFiltration} は次を保持する。
\begin{enumerate}[label=(\roman*),nosep]
  \item 時刻 $n$ ごとの可測空間 $\sigma(n)$。
  \item 単調性 $\sigma(m)\leq \sigma(n)$（$m\leq n$ のとき）。
  \item 包含性 $\sigma(n)\leq \mathcal{M}$（背景可測空間 $\mathcal{M}$ を上界とする）。
  \item 各 $\omega\in\Omega$ に対し、単点集合 $\{\omega\}$ を可測にする最小の時刻が存在。
\end{enumerate}
\end{definition}

この最小時刻を与えるのが \texttt{minLayer} であり、Lean では自然数の選択公理的操作 \texttt{Nat.find} を用いて構成される。これは停止時刻 $\tau(\omega)$ の離散版とみなせる。

\subsection{図式による可視化}

\[
\begin{tikzcd}[column sep=huge,row sep=large]
\sigma(m) \arrow[r,hook,"F.\mathrm{mono}(m\leq n)"] \arrow[d,swap,"\text{包含写像}"] &
\sigma(n) \arrow[d,"\text{包含写像}"] \\
\mathcal{M}(\Omega) \arrow[r,equal] & \mathcal{M}(\Omega)
\end{tikzcd}
\]

図式は、各層 $\sigma(m)$ がより上位の可測空間に包含されること、そして背景可測空間 $\mathcal{M}(\Omega)$ へ安定に埋め込まれることを表す。

\section{構造塔への変換}

Lean の関数 \texttt{DiscreteFiltration.toMeasurableTower} は、フィルトレーションを
\[
T = (\Omega,\ \text{layer}: \mathbb{N} \to \mathrm{MeasurableSpace}\,\Omega,\ \text{minLayer})
\]
という測度論的構造塔に変換する。ここで
\begin{align*}
T.\text{layer}(n) &= \sigma(n),\\
T.\text{minLayer}(\omega) &= \min\{n \mid \{\omega\}\in\sigma(n)\}.
\end{align*}
この変換により、構造塔の理論で用いる単調性や最小性の公理を、そのままフィルトレーションに適用できる。

\section{基本例}

\subsection{自明なフィルトレーション}

\begin{example}
背景可測空間が単点集合を常に可測とする (\texttt{MeasurableSingletonClass}) と仮定すると、Lean の \texttt{trivial} で次が得られる。
\[
\sigma(n) = \mathcal{M}(\Omega)\quad (\forall n),
\]
したがって全ての事象は時刻 $0$ で観測可能である。
\end{example}

\subsection{直積フィルトレーション}

二つのフィルトレーション $F_1,F_2$ の直積構成は
\[
(\sigma_1 \otimes \sigma_2)(n) = \sigma_1(n) \times \sigma_2(n)
\]
として定義される。Lean コード \texttt{product F1 F2} では、既存の \texttt{MeasurableTowerWithMin.prod} を活用して単調性・包含性・単点可測性が保たれることを示している。

\begin{tikzcd}[column sep=huge]
(\Omega_1,\sigma_1(n)) \arrow[r, "F_1.\mathrm{mono}"] \arrow[d,swap,"\times"] &
(\Omega_1,\sigma_1(n+1)) \arrow[d,"\times"] \\
(\Omega_2,\sigma_2(n)) \arrow[r,"F_2.\mathrm{mono}"] &
(\Omega_2,\sigma_2(n+1))
\end{tikzcd}

\section{今後の展望}

\begin{itemize}[nosep]
  \item 連続時間（$\mathbb{R}_{\ge 0}$ 添字）のフィルトレーションへの拡張。
  \item 停止時刻やマルチンゲールを構造塔の射として扱う一般化。
  \item Lean 上での条件付き期待値や Doob 分解の形式化。
\end{itemize}

\section*{謝辞}

本資料および対応する Lean コードは、すべて CODEX (OpenAI) によって自動生成された。luatex での組版を想定しており、必要に応じて \texttt{lualatex} でコンパイルできる。

\end{document}
